\chapter{Physical Parameter Identification and Model Validation}
\label{ch:parameter_identification}

\section{Why Parameter Identification Matters}

You have a DC motor on your bench. The datasheet says its moment of inertia is $J = 0.01 \text{ kg}\cdot\text{m}^2$ and its damping coefficient is $b = 0.1 \text{ N}\cdot\text{m}\cdot\text{s/rad}$. Should you trust these values? The answer is almost always no.

Datasheets give nominal values, often measured under idealized conditions. Your motor is connected to a gearbox, a shaft, a load. The total inertia includes all these components. The damping includes bearing friction, air resistance, and mechanical coupling losses. The real system parameters differ from the datasheet, sometimes by factors of two or more.

This chapter is about measuring the actual physical parameters of your system from experimental data. We will learn how to design experiments, collect data, fit models, validate the results, and quantify uncertainty. This is the bridge between theoretical modeling and real-world control design.

\subsection{When to Trust Datasheets vs Measured Values}

\textbf{Trust datasheets for:}
\begin{itemize}
    \item Electrical parameters that depend only on the component itself (motor resistance, inductance)
    \item Maximum ratings and safety limits (voltage, current, temperature)
    \item Basic physical dimensions and constants
\end{itemize}

\textbf{Measure experimentally for:}
\begin{itemize}
    \item Mechanical parameters involving the complete assembly (inertia, damping)
    \item Coupled dynamics between multiple subsystems
    \item Temperature-dependent parameters in your operating range
    \item Nonlinear characteristics (friction, saturation, deadzone)
    \item Any parameter critical to control performance
\end{itemize}

The rule of thumb: if the parameter appears in your control design equations and affects closed-loop performance, measure it experimentally.

\section{Gray-Box vs Black-Box Identification}

Before we dive into experimental methods, we need to understand two fundamental approaches to system identification.

\subsection{Gray-Box Identification}

In \textbf{gray-box identification}, we know the structure of the governing equations from first principles, but we do not know the numerical values of some parameters. For example, we know a DC motor satisfies:
\begin{equation}
J\ddot{\theta} + b\dot{\theta} = K_t i
\end{equation}
where $J$ is inertia, $b$ is damping, and $K_t$ is the motor torque constant. We know this equation is correct because it follows from Newton's second law and electromagnetic principles. We just need to find the numbers $J$, $b$, and $K_t$ that best match our experimental data.

\textbf{Advantages of gray-box methods:}
\begin{itemize}
    \item Physical parameters have engineering meaning
    \item Models extrapolate better beyond measured operating range
    \item Prior knowledge reduces required data
    \item Parameters can be validated against physical constraints (e.g., $J > 0$)
    \item Easier to diagnose modeling errors
\end{itemize}

\subsection{Black-Box Identification}

In \textbf{black-box identification}, we assume nothing about the internal structure. We simply fit input-output data to a general model form like:
\begin{equation}
G(s) = \frac{b_0 + b_1 s + b_2 s^2 + \cdots}{a_0 + a_1 s + a_2 s^2 + \cdots}
\end{equation}
or a discrete-time difference equation:
\begin{equation}
y[k] = -a_1 y[k-1] - a_2 y[k-2] + b_0 u[k] + b_1 u[k-1]
\end{equation}

The coefficients $a_i$ and $b_i$ are found by minimizing prediction error, but they have no direct physical interpretation.

\textbf{Advantages of black-box methods:}
\begin{itemize}
    \item No prior knowledge of physics required
    \item Can capture complex dynamics we cannot model from first principles
    \item Automated algorithms readily available
    \item Good for control design when physical interpretation is unnecessary
\end{itemize}

\textbf{Disadvantages:}
\begin{itemize}
    \item Models may fit data but violate physics
    \item Poor extrapolation outside measured conditions
    \item Difficult to validate against engineering intuition
    \item Cannot isolate specific physical phenomena
\end{itemize}

\subsection{Which Approach to Use?}

For control engineering, we strongly prefer gray-box methods when feasible. Here is why: if your identified model predicts $J = -0.5 \text{ kg}\cdot\text{m}^2$ (negative inertia), you immediately know something is wrong with the experiment or model structure. A black-box model might fit the data equally well but give you coefficients that look plausible numerically while being physically nonsensical.

Use gray-box identification when:
\begin{itemize}
    \item You understand the governing physics
    \item You need parameters for redesign or optimization
    \item You want to validate manufacturing tolerances
    \item You need to extrapolate to new operating conditions
\end{itemize}

Use black-box identification when:
\begin{itemize}
    \item The physics is too complex to model practically
    \item You only need a model for control design, not physical insight
    \item You are modeling empirical phenomena (e.g., human behavior, market dynamics)
    \item Quick model development is more important than interpretability
\end{itemize}

\section{Experimental Design for Parameter Identification}

The quality of your parameter estimates depends critically on your experimental design. Poor experiments yield unreliable parameters no matter how sophisticated your fitting algorithm.

\subsection{Choosing Excitation Signals}

The input signal you apply to the system must \textbf{excite} all the dynamics you want to identify. This means it must contain frequency content in the bandwidth of interest.

\subsubsection{Step Input}

A step input is the simplest excitation:
\begin{itemize}
    \item Easy to implement
    \item Good for estimating time constants, damping, natural frequency
    \item Excites primarily low-frequency dynamics
    \item Limited information about high-frequency behavior
\end{itemize}

Use step inputs for first-order and second-order systems where transient response characteristics (overshoot, settling time) directly relate to parameters.

\subsubsection{Frequency Sweep (Chirp)}

A frequency sweep varies sinusoidal excitation from low to high frequency:
\begin{equation}
u(t) = A \sin(\omega(t) t), \quad \omega(t) = \omega_{\min} + \frac{\omega_{\max} - \omega_{\min}}{T} t
\end{equation}

\begin{itemize}
    \item Excites all frequencies in a controlled manner
    \item Directly reveals frequency response
    \item Excellent for validating transfer function models
    \item Requires longer experiment time than step test
\end{itemize}

\subsubsection{Pseudorandom Binary Sequence (PRBS)}

A PRBS is a deterministic binary signal that appears random:
\begin{itemize}
    \item Excites broad frequency range with binary (on/off) input
    \item Easy to implement on digital systems
    \item Good for systems with input constraints
    \item Allows statistical analysis of noise
\end{itemize}

\subsection{Practical Considerations}

\subsubsection{Sampling Rate}

Sample at least 10 times faster than the fastest dynamics you want to capture:
\begin{equation}
f_s \geq 10 \cdot f_{\text{bandwidth}}
\end{equation}

For a system with 10 Hz bandwidth, sample at 100 Hz or faster. This ensures you capture the transient response without aliasing.

\subsubsection{Experiment Duration}

Run the experiment long enough to capture the slowest dynamics. For a system with time constant $\tau$, run for at least $5\tau$ to $10\tau$ to see the full settling behavior.

\subsubsection{Signal Amplitude}

The excitation amplitude must be:
\begin{itemize}
    \item Large enough to overcome sensor noise and quantization
    \item Small enough to keep the system in the linear operating region (for linear models)
    \item Within safe operating limits for the hardware
\end{itemize}

A rule of thumb: the signal-to-noise ratio (SNR) should exceed 20 dB (10:1 ratio).

\subsubsection{Initial Conditions}

Start experiments from a known steady-state condition. For many systems, this means starting at rest with zero input. This simplifies analysis because you know $x(0) = 0$ and $\dot{x}(0) = 0$.

\subsubsection{Multiple Trials}

Repeat experiments multiple times to:
\begin{itemize}
    \item Assess repeatability and measurement noise
    \item Average out random disturbances
    \item Detect systematic errors or drift
\end{itemize}

If parameters vary significantly between trials, investigate environmental factors (temperature, backlash, wear).

\section{Least-Squares Parameter Estimation}

Least-squares estimation is the workhorse of parameter identification. We will develop it carefully from first principles, including the geometric interpretation that makes it intuitive.

\subsection{The Problem Statement}

Suppose we have a linear model:
\begin{equation}
y = \theta_1 \phi_1 + \theta_2 \phi_2 + \cdots + \theta_n \phi_n
\end{equation}
where $y$ is the measured output, $\phi_i$ are known functions of time or past measurements (called \textbf{regressors}), and $\theta_i$ are unknown parameters we want to estimate.

We collect $N$ data points at times $t_1, t_2, \ldots, t_N$. For each measurement $k$, we have:
\begin{equation}
y[k] = \theta_1 \phi_1[k] + \theta_2 \phi_2[k] + \cdots + \theta_n \phi_n[k] + e[k]
\end{equation}
where $e[k]$ is measurement noise and modeling error.

In matrix form:
\begin{equation}
\mathbf{y} = \mathbf{\Phi} \boldsymbol{\theta} + \mathbf{e}
\end{equation}
where:
\begin{equation}
\mathbf{y} = \begin{bmatrix} y[1] \\ y[2] \\ \vdots \\ y[N] \end{bmatrix}, \quad
\mathbf{\Phi} = \begin{bmatrix}
\phi_1[1] & \phi_2[1] & \cdots & \phi_n[1] \\
\phi_1[2] & \phi_2[2] & \cdots & \phi_n[2] \\
\vdots & \vdots & \ddots & \vdots \\
\phi_1[N] & \phi_2[N] & \cdots & \phi_n[N]
\end{bmatrix}, \quad
\boldsymbol{\theta} = \begin{bmatrix} \theta_1 \\ \theta_2 \\ \vdots \\ \theta_n \end{bmatrix}
\end{equation}

The matrix $\mathbf{\Phi}$ is called the \textbf{regressor matrix} or \textbf{design matrix}.

Our goal is to find $\boldsymbol{\theta}$ that minimizes the sum of squared errors:
\begin{equation}
J(\boldsymbol{\theta}) = \sum_{k=1}^N e[k]^2 = \sum_{k=1}^N \left( y[k] - \sum_{i=1}^n \theta_i \phi_i[k] \right)^2
\end{equation}

In matrix notation:
\begin{equation}
J(\boldsymbol{\theta}) = (\mathbf{y} - \mathbf{\Phi}\boldsymbol{\theta})^T (\mathbf{y} - \mathbf{\Phi}\boldsymbol{\theta}) = \|\mathbf{y} - \mathbf{\Phi}\boldsymbol{\theta}\|^2
\end{equation}

\subsection{Geometric Interpretation}

Here is the key insight: $\mathbf{y}$ is a vector in $\mathbb{R}^N$ representing our measurements. $\mathbf{\Phi}\boldsymbol{\theta}$ is also a vector in $\mathbb{R}^N$ representing our model predictions. But $\mathbf{\Phi}\boldsymbol{\theta}$ is constrained to lie in the column space of $\mathbf{\Phi}$, which is an $n$-dimensional subspace of $\mathbb{R}^N$ (assuming $\mathbf{\Phi}$ has full rank).

Least-squares finds the point in this subspace closest to $\mathbf{y}$. This is a projection problem.

\begin{figure}[h]
\centering
\begin{tikzpicture}[scale=1.5]
    % Coordinate system
    \draw[->] (0,0) -- (4,0) node[right] {$\phi_1$};
    \draw[->] (0,0) -- (0,4) node[above] {$\phi_2$};

    % Plane spanned by columns of Phi
    \fill[Atlantic!20] (0,0) -- (3.5,0) -- (3.5,3.5) -- (0,3.5) -- cycle;
    \node at (3,3.2) {Column space of $\mathbf{\Phi}$};

    % Data vector y
    \draw[->,thick,Garnet] (0,0) -- (2.5,3.5) node[above] {$\mathbf{y}$ (data)};

    % Projection Phi*theta_hat
    \draw[->,thick,Congaree] (0,0) -- (2.5,2.5) node[right] {$\mathbf{\Phi}\hat{\boldsymbol{\theta}}$ (model)};

    % Residual
    \draw[->,dashed,thick,Honeycomb] (2.5,2.5) -- (2.5,3.5) node[midway,right] {$\mathbf{e}$ (residual)};

    % Right angle
    \draw (2.3,2.5) -- (2.3,2.7) -- (2.5,2.7);
\end{tikzpicture}
\caption{Geometric interpretation of least-squares estimation. The data vector $\mathbf{y}$ is projected orthogonally onto the column space of $\mathbf{\Phi}$. The optimal parameter vector $\hat{\boldsymbol{\theta}}$ yields the projection $\mathbf{\Phi}\hat{\boldsymbol{\theta}}$ closest to $\mathbf{y}$.}
\label{fig:least_squares_geometry}
\end{figure}

The residual vector $\mathbf{e} = \mathbf{y} - \mathbf{\Phi}\hat{\boldsymbol{\theta}}$ is orthogonal to the column space of $\mathbf{\Phi}$. This orthogonality condition gives us the solution.

\subsection{Derivation of the Least-Squares Solution}

To minimize $J(\boldsymbol{\theta})$, we take the derivative with respect to $\boldsymbol{\theta}$ and set it to zero.

Start with:
\begin{equation}
J(\boldsymbol{\theta}) = (\mathbf{y} - \mathbf{\Phi}\boldsymbol{\theta})^T (\mathbf{y} - \mathbf{\Phi}\boldsymbol{\theta})
\end{equation}

Expand:
\begin{equation}
J(\boldsymbol{\theta}) = \mathbf{y}^T\mathbf{y} - \mathbf{y}^T\mathbf{\Phi}\boldsymbol{\theta} - \boldsymbol{\theta}^T\mathbf{\Phi}^T\mathbf{y} + \boldsymbol{\theta}^T\mathbf{\Phi}^T\mathbf{\Phi}\boldsymbol{\theta}
\end{equation}

Since $\mathbf{y}^T\mathbf{\Phi}\boldsymbol{\theta}$ is a scalar, it equals its transpose:
\begin{equation}
\mathbf{y}^T\mathbf{\Phi}\boldsymbol{\theta} = (\mathbf{y}^T\mathbf{\Phi}\boldsymbol{\theta})^T = \boldsymbol{\theta}^T\mathbf{\Phi}^T\mathbf{y}
\end{equation}

Therefore:
\begin{equation}
J(\boldsymbol{\theta}) = \mathbf{y}^T\mathbf{y} - 2\boldsymbol{\theta}^T\mathbf{\Phi}^T\mathbf{y} + \boldsymbol{\theta}^T\mathbf{\Phi}^T\mathbf{\Phi}\boldsymbol{\theta}
\end{equation}

Take the derivative with respect to $\boldsymbol{\theta}$:
\begin{equation}
\frac{\partial J}{\partial \boldsymbol{\theta}} = -2\mathbf{\Phi}^T\mathbf{y} + 2\mathbf{\Phi}^T\mathbf{\Phi}\boldsymbol{\theta}
\end{equation}

Set equal to zero:
\begin{equation}
-2\mathbf{\Phi}^T\mathbf{y} + 2\mathbf{\Phi}^T\mathbf{\Phi}\boldsymbol{\theta} = 0
\end{equation}

Simplify:
\begin{equation}
\mathbf{\Phi}^T\mathbf{\Phi}\boldsymbol{\theta} = \mathbf{\Phi}^T\mathbf{y}
\end{equation}

These are called the \textbf{normal equations}. If $\mathbf{\Phi}^T\mathbf{\Phi}$ is invertible (which requires $\mathbf{\Phi}$ to have full column rank), we solve for $\boldsymbol{\theta}$:

\begin{equation}
\boxed{\hat{\boldsymbol{\theta}} = (\mathbf{\Phi}^T\mathbf{\Phi})^{-1}\mathbf{\Phi}^T\mathbf{y}}
\end{equation}

This is the \textbf{least-squares solution}. The matrix $(\mathbf{\Phi}^T\mathbf{\Phi})^{-1}\mathbf{\Phi}^T$ is called the \textbf{pseudo-inverse} of $\mathbf{\Phi}$.

\subsection{When Does Least-Squares Work?}

The least-squares solution exists and is unique if and only if $\mathbf{\Phi}$ has full column rank, meaning the columns of $\mathbf{\Phi}$ are linearly independent. Physically, this means:

\begin{itemize}
    \item Each parameter affects the output in a distinguishable way
    \item Your experiment provides enough independent information
    \item No parameter is redundant or perfectly correlated with others
\end{itemize}

\textbf{Common causes of rank deficiency:}
\begin{enumerate}
    \item Insufficient excitation: the input does not excite all modes
    \item Redundant parameterization: two parameters always appear together
    \item Poor experiment design: operating at a single frequency or steady state
\end{enumerate}

\subsection{Weighted Least-Squares}

When measurements have different noise levels, we can weight them accordingly. The weighted least-squares criterion is:
\begin{equation}
J(\boldsymbol{\theta}) = \sum_{k=1}^N w_k e[k]^2 = (\mathbf{y} - \mathbf{\Phi}\boldsymbol{\theta})^T \mathbf{W} (\mathbf{y} - \mathbf{\Phi}\boldsymbol{\theta})
\end{equation}
where $\mathbf{W} = \text{diag}(w_1, w_2, \ldots, w_N)$ is a diagonal matrix of weights.

The solution becomes:
\begin{equation}
\hat{\boldsymbol{\theta}} = (\mathbf{\Phi}^T\mathbf{W}\mathbf{\Phi})^{-1}\mathbf{\Phi}^T\mathbf{W}\mathbf{y}
\end{equation}

Set $w_k = 1/\sigma_k^2$ where $\sigma_k^2$ is the variance of measurement $k$. This gives more weight to accurate measurements and less to noisy ones.

\section{Parameter Estimation from Step Response}

Let us apply least-squares to a concrete problem: identifying the inertia and damping of a DC motor from a step response.

\subsection{First-Order System: DC Motor Velocity}

Consider a DC motor under armature voltage control, where we measure angular velocity $\omega(t)$. Neglecting inductance, the dynamics are:
\begin{equation}
J\dot{\omega} + b\omega = K_t i = \frac{K_t}{R_a}(V - K_e\omega)
\end{equation}

Rearranging:
\begin{equation}
\dot{\omega} = -\frac{b}{J}\omega - \frac{K_tK_e}{JR_a}\omega + \frac{K_t}{JR_a}V = -\frac{1}{\tau}\omega + \frac{K}{\tau}V
\end{equation}

where:
\begin{align}
\tau &= \frac{JR_a}{bR_a + K_tK_e} \quad \text{(time constant)} \\
K &= \frac{K_t}{bR_a + K_tK_e} \quad \text{(DC gain)}
\end{align}

For a step input $V(t) = V_0$ starting from rest, the analytical solution is:
\begin{equation}
\omega(t) = KV_0(1 - e^{-t/\tau})
\end{equation}

From this, we can identify $\tau$ and $K$ by curve fitting. Once we have $\tau$ and $K$, we can extract the physical parameters $J$ and $b$ if we know the electrical constants $R_a$, $K_t$, and $K_e$ from the motor datasheet.

\subsection{Method 1: Analytical Curve Fitting}

Given measured data $\{\omega[k], t_k\}$ for $k = 1, 2, \ldots, N$, we fit:
\begin{equation}
\omega[k] = \theta_1(1 - e^{-t_k/\theta_2}) + e[k]
\end{equation}
where $\theta_1 = KV_0$ and $\theta_2 = \tau$ are the parameters to estimate.

This is a nonlinear least-squares problem because $\tau$ appears in the exponential. We can solve it using:
\begin{itemize}
    \item Nonlinear optimization (Gauss-Newton, Levenberg-Marquardt)
    \item Logarithmic transformation (less accurate but analytical)
    \item Graphical methods (time to 63.2\% of final value gives $\tau$)
\end{itemize}

\textbf{Graphical Method (Quick and Dirty):}
\begin{enumerate}
    \item Find the steady-state value: $\omega_{\infty} \approx \omega[N]$
    \item Compute $0.632 \times \omega_{\infty}$
    \item Find the time $t_{0.632}$ when $\omega(t)$ crosses this value
    \item Estimate $\tau \approx t_{0.632}$
    \item Estimate $K = \omega_{\infty}/V_0$
\end{enumerate}

This works well for clean data but is sensitive to noise.

\subsection{Method 2: Differential Equation Form}

A more robust approach uses the differential equation directly. Discretize using backward differences:
\begin{equation}
\dot{\omega}[k] \approx \frac{\omega[k] - \omega[k-1]}{\Delta t}
\end{equation}

Substitute into the continuous-time model:
\begin{equation}
\frac{\omega[k] - \omega[k-1]}{\Delta t} = -\frac{1}{\tau}\omega[k] + \frac{K}{\tau}V[k]
\end{equation}

Rearrange:
\begin{equation}
\omega[k] = \left(1 - \frac{\Delta t}{\tau}\right)\omega[k-1] + \frac{K\Delta t}{\tau}V[k]
\end{equation}

Define:
\begin{align}
\theta_1 &= 1 - \frac{\Delta t}{\tau} \\
\theta_2 &= \frac{K\Delta t}{\tau}
\end{align}

This gives a linear regression:
\begin{equation}
\omega[k] = \theta_1 \omega[k-1] + \theta_2 V[k] + e[k]
\end{equation}

The regressor matrix is:
\begin{equation}
\mathbf{\Phi} = \begin{bmatrix}
\omega[1] & V[2] \\
\omega[2] & V[3] \\
\vdots & \vdots \\
\omega[N-1] & V[N]
\end{bmatrix}, \quad
\mathbf{y} = \begin{bmatrix}
\omega[2] \\
\omega[3] \\
\vdots \\
\omega[N]
\end{bmatrix}
\end{equation}

Solve for $\hat{\boldsymbol{\theta}} = (\mathbf{\Phi}^T\mathbf{\Phi})^{-1}\mathbf{\Phi}^T\mathbf{y}$.

Recover $\tau$ and $K$ from:
\begin{align}
\tau &= \frac{\Delta t}{1 - \theta_1} \\
K &= \frac{\theta_2(1 - \theta_1)}{\Delta t} = \theta_2 / \tau
\end{align}

\textbf{Advantages of this method:}
\begin{itemize}
    \item Linear in parameters (no nonlinear optimization needed)
    \item Naturally handles noisy data through least-squares averaging
    \item Works even if the system does not reach steady state
    \item Easily extends to higher-order systems
\end{itemize}

\subsection{Second-Order System: Motor with Inertial Load}

Now consider a motor driving a load with position feedback. The dynamics are:
\begin{equation}
J\ddot{\theta} + b\dot{\theta} = K_t i
\end{equation}

With voltage control and back-EMF, this becomes a second-order system. In transfer function form:
\begin{equation}
G(s) = \frac{\Theta(s)}{V(s)} = \frac{K}{\tau^2 s^2 + 2\zeta\tau s + 1}
\end{equation}

where:
\begin{align}
\tau &= \sqrt{\frac{JR_a}{bR_a + K_tK_e}} \quad \text{(natural period)} \\
\zeta &= \frac{b}{2}\sqrt{\frac{R_a}{J(bR_a + K_tK_e)}} \quad \text{(damping ratio)}
\end{align}

From a step response, we can measure:
\begin{itemize}
    \item Rise time $t_r$
    \item Peak time $t_p$
    \item Percent overshoot $PO$
    \item Settling time $t_s$
\end{itemize}

The damping ratio relates to percent overshoot:
\begin{equation}
PO = e^{-\pi\zeta/\sqrt{1-\zeta^2}} \times 100\%
\end{equation}

Solving for $\zeta$:
\begin{equation}
\zeta = \frac{-\ln(PO/100)}{\sqrt{\pi^2 + \ln^2(PO/100)}}
\end{equation}

The natural frequency $\omega_n = 1/\tau$ relates to peak time:
\begin{equation}
t_p = \frac{\pi}{\omega_n\sqrt{1-\zeta^2}} \quad \Rightarrow \quad \omega_n = \frac{\pi}{t_p\sqrt{1-\zeta^2}}
\end{equation}

Once we have $\zeta$ and $\omega_n$, we can extract $J$ and $b$ if we know the electrical parameters.

\textbf{Procedure:}
\begin{enumerate}
    \item Apply a step input and record position $\theta(t)$
    \item Identify the final value $\theta_{\infty}$
    \item Measure the first peak value $\theta_{\max}$
    \item Calculate $PO = (\theta_{\max} - \theta_{\infty})/\theta_{\infty} \times 100\%$
    \item Measure the time to peak $t_p$
    \item Compute $\zeta$ from overshoot
    \item Compute $\omega_n$ from peak time
    \item Extract $J$ and $b$ using motor constants
\end{enumerate}

\section{Parameter Estimation from Frequency Response}

Frequency response methods provide an alternative approach, especially useful when time-domain transients are difficult to measure or when the system operates continuously.

\subsection{Experimental Procedure}

\begin{enumerate}
    \item Apply a sinusoidal input at frequency $\omega$: $u(t) = A\sin(\omega t)$
    \item Wait for transients to decay (typically 5-10 periods)
    \item Measure the steady-state output: $y(t) = B\sin(\omega t + \phi)$
    \item Compute magnitude ratio: $|G(j\omega)| = B/A$
    \item Compute phase shift: $\angle G(j\omega) = \phi$
    \item Repeat for multiple frequencies spanning the bandwidth
    \item Fit a transfer function model to the frequency response data
\end{enumerate}

\subsection{First-Order System Frequency Response}

For the first-order model:
\begin{equation}
G(s) = \frac{K}{\tau s + 1}
\end{equation}

The frequency response is:
\begin{equation}
G(j\omega) = \frac{K}{\tau j\omega + 1} = \frac{K}{1 + j\omega\tau}
\end{equation}

Magnitude:
\begin{equation}
|G(j\omega)| = \frac{K}{\sqrt{1 + (\omega\tau)^2}}
\end{equation}

Phase:
\begin{equation}
\angle G(j\omega) = -\arctan(\omega\tau)
\end{equation}

\textbf{Parameter Extraction:}
\begin{itemize}
    \item At $\omega = 0$: $|G(0)| = K$ (DC gain)
    \item At $\omega = 1/\tau$: $|G(j/\tau)| = K/\sqrt{2} = 0.707K$ (this is the \textbf{cutoff frequency})
    \item At $\omega = 1/\tau$: $\angle G(j/\tau) = -45°$
\end{itemize}

So we can identify $\tau$ by finding the frequency where the phase is $-45°$ or the magnitude drops to $0.707$ of the DC gain.

\subsection{Least-Squares Fitting to Frequency Response}

Given frequency response measurements $\{|G(j\omega_k)|, \angle G(j\omega_k)\}$ for $k = 1, 2, \ldots, N$, we want to fit a model. For a first-order system, define:
\begin{equation}
|G(j\omega)|^2 = \frac{K^2}{1 + (\omega\tau)^2}
\end{equation}

Rearrange:
\begin{equation}
|G(j\omega)|^2 (1 + (\omega\tau)^2) = K^2
\end{equation}

\begin{equation}
|G(j\omega)|^2 = K^2 - |G(j\omega)|^2(\omega\tau)^2
\end{equation}

Define $\theta_1 = K^2$ and $\theta_2 = -K^2\tau^2$:
\begin{equation}
|G(j\omega_k)|^2 = \theta_1 + \theta_2 \omega_k^2
\end{equation}

This is linear in $\theta_1$ and $\theta_2$. The regressor matrix is:
\begin{equation}
\mathbf{\Phi} = \begin{bmatrix}
1 & \omega_1^2 \\
1 & \omega_2^2 \\
\vdots & \vdots \\
1 & \omega_N^2
\end{bmatrix}, \quad
\mathbf{y} = \begin{bmatrix}
|G(j\omega_1)|^2 \\
|G(j\omega_2)|^2 \\
\vdots \\
|G(j\omega_N)|^2
\end{bmatrix}
\end{equation}

Solve $\hat{\boldsymbol{\theta}} = (\mathbf{\Phi}^T\mathbf{\Phi})^{-1}\mathbf{\Phi}^T\mathbf{y}$.

Extract:
\begin{align}
K &= \sqrt{\theta_1} \\
\tau &= \sqrt{-\theta_2/\theta_1}
\end{align}

For higher-order systems, use nonlinear optimization to fit magnitude and phase simultaneously.

\section{Practical Challenges in Parameter Identification}

Real experiments are messy. Here are the most common challenges and how to handle them.

\subsection{Measurement Noise}

All sensors have noise. High-frequency noise corrupts derivative estimates, which are essential for identifying dynamics.

\textbf{Solutions:}
\begin{itemize}
    \item Low-pass filter the data before differentiation
    \item Use numerical differentiation schemes that smooth noise (e.g., Savitzky-Golay filter)
    \item Increase signal amplitude to improve signal-to-noise ratio
    \item Average multiple experimental runs
    \item Use integral-based identification methods that naturally smooth noise
\end{itemize}

\subsection{Unmeasured Disturbances}

Disturbances you cannot measure (e.g., air currents, thermal drift, electromagnetic interference) corrupt the data.

\textbf{Solutions:}
\begin{itemize}
    \item Isolate the system mechanically and electrically
    \item Run experiments quickly to minimize drift
    \item Use differential measurements (measure relative quantities)
    \item Model disturbances as additional states and estimate them
\end{itemize}

\subsection{Nonlinearity}

Real systems are nonlinear. Friction is not proportional to velocity. Springs have nonlinear stiffness. Actuators saturate.

\textbf{Approaches:}
\begin{enumerate}
    \item \textbf{Linearize around an operating point:} Use small excitations to stay in the quasi-linear region. Identify separate models for different operating points and use gain scheduling.

    \item \textbf{Identify nonlinear models:} Fit polynomial or piecewise-linear models. For example, Coulomb friction:
    \begin{equation}
    f(\dot{x}) = b\dot{x} + f_c \text{sign}(\dot{x})
    \end{equation}
    This adds one more parameter ($f_c$) to estimate.

    \item \textbf{Use gray-box nonlinear structures:} If you know the form of the nonlinearity (e.g., saturation, deadzone), fit those specific functions.
\end{enumerate}

\subsection{Coupling and Cross-Talk}

In MIMO systems, inputs affect multiple outputs and vice versa. Identifying one parameter changes the fit for another.

\textbf{Solutions:}
\begin{itemize}
    \item Design experiments to excite one input at a time when possible
    \item Use MIMO identification methods that account for coupling
    \item Iterate: identify one subsystem, fix those parameters, identify the next
\end{itemize}

\subsection{Insufficient Excitation}

If your input does not excite certain modes, you cannot identify the corresponding parameters.

\textbf{Example:} Operating a motor at constant speed does not reveal inertia because $\ddot{\theta} = 0$.

\textbf{Solutions:}
\begin{itemize}
    \item Use rich excitation signals (step, chirp, PRBS)
    \item Check the rank of $\mathbf{\Phi}^T\mathbf{\Phi}$ numerically (condition number)
    \item Design experiments specifically to excite unobservable modes
\end{itemize}

\section{Model Validation Techniques}

Identifying parameters is only half the battle. We must validate that the model accurately represents the real system.

\subsection{Residual Analysis}

The residual is the difference between measured and predicted outputs:
\begin{equation}
e[k] = y[k] - \hat{y}[k]
\end{equation}

For a good model with Gaussian noise, the residuals should:
\begin{itemize}
    \item Have zero mean: $\bar{e} \approx 0$
    \item Be uncorrelated (white noise): $\text{corr}(e[k], e[k+\tau]) \approx 0$ for $\tau \neq 0$
    \item Be normally distributed
    \item Show no systematic patterns
\end{itemize}

\textbf{Diagnostic Plots:}
\begin{enumerate}
    \item \textbf{Residual vs time:} Should look like random scatter with no trends
    \item \textbf{Residual autocorrelation:} Should be zero except at lag zero
    \item \textbf{Residual histogram:} Should be approximately Gaussian
    \item \textbf{Residual vs predicted output:} Should show no correlation (constant variance)
\end{enumerate}

If residuals show patterns, the model structure is inadequate. Consider:
\begin{itemize}
    \item Adding more parameters (higher-order model)
    \item Including nonlinear terms
    \item Modeling disturbances explicitly
\end{itemize}

\subsection{Cross-Validation}

Never validate on the same data used for identification. The procedure is:
\begin{enumerate}
    \item Split data into training set (70-80\%) and validation set (20-30\%)
    \item Identify parameters using training set
    \item Predict outputs on validation set
    \item Compute validation error metrics
\end{enumerate}

If validation error is significantly larger than training error, the model is \textbf{overfitting}.

\subsection{Physical Consistency Checks}

Before accepting identified parameters, verify they make physical sense:
\begin{itemize}
    \item Inertia $J > 0$
    \item Damping $b \geq 0$ (dissipative systems)
    \item Stiffness $k > 0$
    \item Capacitance $C > 0$, inductance $L > 0$, resistance $R > 0$
    \item Time constants match expected order of magnitude
\end{itemize}

If parameters violate physics, there is likely:
\begin{itemize}
    \item An error in the model structure
    \item Insufficient data quality
    \item Numerical ill-conditioning
\end{itemize}

\subsection{Frequency Response Validation}

Compare the frequency response of the identified model to experimental measurements:
\begin{itemize}
    \item Plot Bode magnitude and phase
    \item Check that resonances align
    \item Verify DC gain matches
    \item Confirm high-frequency roll-off rate
\end{itemize}

Frequency-domain errors reveal specific deficiencies:
\begin{itemize}
    \item DC gain mismatch: wrong steady-state parameters
    \item Incorrect resonant frequency: wrong natural frequency or inertia
    \item Wrong phase slope: unmodeled time delay or additional poles/zeros
\end{itemize}

\subsection{Step Response Validation}

Apply a step input not used during identification and compare the model prediction to the actual response:
\begin{itemize}
    \item Rise time should match
    \item Overshoot should be similar
    \item Settling time should align
    \item Steady-state value should be correct
\end{itemize}

This is the ultimate test: does the model predict behavior on new data?

\section{Uncertainty Quantification}

Parameter estimates are uncertain due to measurement noise and finite data. We need to quantify this uncertainty to assess model reliability.

\subsection{Parameter Covariance Matrix}

Under the assumption of Gaussian noise with variance $\sigma^2$, the covariance of the parameter estimates is:
\begin{equation}
\text{Cov}(\hat{\boldsymbol{\theta}}) = \sigma^2 (\mathbf{\Phi}^T\mathbf{\Phi})^{-1}
\end{equation}

Estimate the noise variance from the residuals:
\begin{equation}
\hat{\sigma}^2 = \frac{1}{N - n} \sum_{k=1}^N e[k]^2
\end{equation}

where $n$ is the number of parameters.

The diagonal elements of $\text{Cov}(\hat{\boldsymbol{\theta}})$ give the variance of each parameter estimate:
\begin{equation}
\text{Var}(\hat{\theta}_i) = [\text{Cov}(\hat{\boldsymbol{\theta}})]_{ii}
\end{equation}

The standard error is:
\begin{equation}
\text{SE}(\hat{\theta}_i) = \sqrt{\text{Var}(\hat{\theta}_i)}
\end{equation}

A 95\% confidence interval for $\theta_i$ is approximately:
\begin{equation}
\hat{\theta}_i \pm 2 \cdot \text{SE}(\hat{\theta}_i)
\end{equation}

\subsection{Interpreting Uncertainty}

\textbf{Large standard errors indicate:}
\begin{itemize}
    \item High measurement noise
    \item Insufficient data
    \item Poor experimental design (low signal-to-noise ratio)
    \item Ill-conditioned regressor matrix (parameters are correlated)
\end{itemize}

\textbf{How to reduce uncertainty:}
\begin{itemize}
    \item Collect more data
    \item Increase excitation amplitude
    \item Use better sensors
    \item Redesign the experiment to decorrelate parameters
\end{itemize}

\subsection{Monte Carlo Uncertainty Propagation}

To understand how parameter uncertainty affects model predictions:
\begin{enumerate}
    \item Generate $M$ samples of parameters from the covariance: $\boldsymbol{\theta}^{(i)} \sim \mathcal{N}(\hat{\boldsymbol{\theta}}, \text{Cov}(\hat{\boldsymbol{\theta}}))$
    \item Simulate the model with each parameter set
    \item Compute statistics (mean, variance, confidence bounds) of the predictions
\end{enumerate}

This gives you prediction intervals that account for parameter uncertainty.

\section{When to Trust Identified Parameters vs Nominal Values}

This is the engineering judgment question. Here are guidelines:

\subsection{Trust Identified Parameters When:}
\begin{itemize}
    \item Standard errors are small (less than 10\% of the parameter value)
    \item Validation error is acceptable on new data
    \item Parameters are physically plausible and consistent with expectations
    \item Multiple experiments give consistent results
    \item Residuals show no systematic structure
\end{itemize}

\subsection{Be Skeptical When:}
\begin{itemize}
    \item Parameters differ wildly from nominal values (factors of 5-10)
    \item Standard errors are large (50\% or more of parameter value)
    \item Validation error is much larger than training error
    \item Parameters violate physical constraints
    \item Residuals show clear patterns
\end{itemize}

\subsection{Use a Hybrid Approach:}
\begin{itemize}
    \item Fix well-known parameters (e.g., motor torque constant from datasheet)
    \item Identify only uncertain parameters (e.g., total inertia including load)
    \item Use nominal values as initial guesses in optimization
    \item Regularize: penalize deviations from nominal values
\end{itemize}

\section{Iterative Refinement: Model → Test → Refine → Test}

Parameter identification is rarely one-shot. The process is iterative:

\subsection{Step 1: Initial Model from First Principles}
Build a preliminary model using nominal parameters and physics-based equations.

\subsection{Step 2: Design Experiment}
Choose excitation signals, sampling rates, and measurement points based on the initial model's predicted dynamics.

\subsection{Step 3: Collect Data}
Run the experiment carefully, monitoring data quality in real time.

\subsection{Step 4: Identify Parameters}
Fit the model to data using least-squares or other estimation methods.

\subsection{Step 5: Validate}
Check residuals, validate on new data, verify physical plausibility.

\subsection{Step 6: Refine Model Structure}
If validation fails:
\begin{itemize}
    \item Add unmodeled dynamics (friction, flexibility, time delays)
    \item Include nonlinearities
    \item Increase model order
\end{itemize}

\subsection{Step 7: Iterate}
Repeat steps 4-6 until validation is satisfactory.

\subsection{Step 8: Deploy}
Use the validated model for control design.

\subsection{Step 9: Monitor in Operation}
Track control performance. If it degrades, revisit identification (parameters may drift due to wear, temperature, aging).

\section{Example 1: DC Motor Inertia and Damping from Step Test}

\subsection{System Description}

We have a DC motor (Maxon EC45 flat) driving a disk load through a rigid shaft. We want to identify the total inertia $J$ and damping $b$ of the rotating system.

\textbf{Known parameters from datasheet:}
\begin{itemize}
    \item Armature resistance: $R_a = 0.75 \, \Omega$
    \item Torque constant: $K_t = 0.0265 \, \text{Nm/A}$
    \item Back-EMF constant: $K_e = 0.0265 \, \text{V}\cdot\text{s/rad}$
    \item Rotor inertia (motor only): $J_{\text{motor}} = 25 \times 10^{-6} \, \text{kg}\cdot\text{m}^2$
\end{itemize}

\textbf{Additional load:} Disk with diameter 10 cm, mass 0.2 kg.

\textbf{Expected total inertia:}
\begin{equation}
J_{\text{expected}} = J_{\text{motor}} + \frac{1}{2}mr^2 = 25\times10^{-6} + \frac{1}{2}(0.2)(0.05)^2 = 275 \times 10^{-6} \, \text{kg}\cdot\text{m}^2
\end{equation}

\subsection{Experimental Setup}

\begin{itemize}
    \item Encoder: 2000 counts per revolution (CPR), sampled at 1 kHz
    \item Controller: Arduino with motor driver
    \item Input: Step voltage from 0 to 5 V
    \item Measurement: Angular velocity $\omega(t)$ computed by differentiating encoder counts
\end{itemize}

\subsection{Data Collection}

We apply a 5 V step and record velocity for 2 seconds (2000 samples). The motor accelerates from rest and reaches approximately 800 rad/s steady state.

\subsection{Method 1: Graphical Time Constant}

From the plot, we observe:
\begin{itemize}
    \item Steady-state velocity: $\omega_{\infty} = 798 \, \text{rad/s}$
    \item 63.2\% of final value: $0.632 \times 798 = 504 \, \text{rad/s}$
    \item Time to reach 504 rad/s: $t_{0.632} = 0.042 \, \text{s}$
\end{itemize}

Thus:
\begin{equation}
\tau \approx 0.042 \, \text{s}
\end{equation}

The DC gain is:
\begin{equation}
K = \frac{\omega_{\infty}}{V_0} = \frac{798}{5} = 159.6 \, \text{rad/s/V}
\end{equation}

\subsection{Method 2: Least-Squares Differential Equation}

Discretize the model:
\begin{equation}
\omega[k] = \theta_1 \omega[k-1] + \theta_2 V[k]
\end{equation}

where:
\begin{align}
\theta_1 &= 1 - \frac{\Delta t}{\tau} = 1 - \frac{0.001}{\tau} \\
\theta_2 &= \frac{K\Delta t}{\tau}
\end{align}

Construct regressor matrix with $N = 2000$ samples:
\begin{equation}
\mathbf{\Phi} = \begin{bmatrix}
\omega[1] & 5 \\
\omega[2] & 5 \\
\vdots & \vdots \\
\omega[1999] & 5
\end{bmatrix}, \quad
\mathbf{y} = \begin{bmatrix}
\omega[2] \\
\omega[3] \\
\vdots \\
\omega[2000]
\end{bmatrix}
\end{equation}

Compute:
\begin{equation}
\hat{\boldsymbol{\theta}} = (\mathbf{\Phi}^T\mathbf{\Phi})^{-1}\mathbf{\Phi}^T\mathbf{y}
\end{equation}

\textbf{Result:}
\begin{equation}
\hat{\boldsymbol{\theta}} = \begin{bmatrix} 0.9765 \\ 3.745 \end{bmatrix}
\end{equation}

Extract:
\begin{align}
\tau &= \frac{\Delta t}{1 - \theta_1} = \frac{0.001}{1 - 0.9765} = 0.0426 \, \text{s} \\
K &= \frac{\theta_2}{\Delta t} \tau = \frac{3.745}{0.001} \times 0.0426 = 159.5 \, \text{rad/s/V}
\end{align}

This matches the graphical method well.

\subsection{Extracting Physical Parameters}

The theoretical relationships are:
\begin{align}
\tau &= \frac{JR_a}{bR_a + K_tK_e} \\
K &= \frac{K_t}{bR_a + K_tK_e}
\end{align}

From $K$:
\begin{equation}
bR_a + K_tK_e = \frac{K_t}{K} = \frac{0.0265}{159.5} = 1.662 \times 10^{-4}
\end{equation}

From $\tau$:
\begin{equation}
J = \frac{\tau(bR_a + K_tK_e)}{R_a} = \frac{0.0426 \times 1.662 \times 10^{-4}}{0.75} = 9.44 \times 10^{-6} \, \text{kg}\cdot\text{m}^2
\end{equation}

\textbf{Problem:} This is much smaller than expected! We expected $275 \times 10^{-6}$ but got $9.44 \times 10^{-6}$.

\subsection{Diagnosis}

The issue is that we are identifying the \textbf{electrical time constant}, not the mechanical one. When inductance $L$ is neglected, the model assumes infinite current rise time. In reality, the armature inductance causes current (and thus torque) to build gradually.

For a more accurate mechanical parameter identification, we need to:
\begin{enumerate}
    \item Include inductance in the model (second-order electrical dynamics)
    \item Use a slower time scale (measure position, not velocity)
    \item Apply a current-controlled input (bypassing electrical dynamics)
\end{enumerate}

\subsection{Refined Experiment: Position Step Response}

Instead of velocity, we measure position $\theta(t)$ in response to a voltage step.

The transfer function from voltage to position is:
\begin{equation}
\frac{\Theta(s)}{V(s)} = \frac{K_t/R_a}{Js^2 + (b + K_tK_e/R_a)s}
\end{equation}

This is a second-order system (double integrator with damping). The step response will show:
\begin{itemize}
    \item Initial acceleration phase
    \item Linear ramp once steady-state velocity is reached
\end{itemize}

From the position data, we can fit:
\begin{equation}
\theta(t) = \frac{1}{2}at^2 + v_0 t + \theta_0
\end{equation}

The acceleration $a = \tau_{\text{motor}}/J$ where $\tau_{\text{motor}}$ is the motor torque.

\textbf{Measured position data:} After 1 second, $\theta = 35 \, \text{rad}$.

The motor torque is:
\begin{equation}
\tau_{\text{motor}} = K_t i_{\infty} = K_t \frac{V}{R_a} = 0.0265 \times \frac{5}{0.75} = 0.1767 \, \text{Nm}
\end{equation}

At steady state (ignoring damping, which is small):
\begin{equation}
\omega_{\infty} = \frac{K_t V}{K_e K_t + bR_a} \approx \frac{V}{K_e} = \frac{5}{0.0265} = 188.7 \, \text{rad/s}
\end{equation}

Wait, this does not match our measured 798 rad/s. Let me recalculate.

Actually, the steady-state velocity is:
\begin{equation}
\omega_{\infty} = \frac{K_t/R_a}{K_e K_t / R_a + b} V
\end{equation}

If $b$ is very small:
\begin{equation}
\omega_{\infty} \approx \frac{V}{K_e} = \frac{5}{0.0265} = 188.7 \, \text{rad/s}
\end{equation}

But we measured 798 rad/s. This suggests either:
\begin{itemize}
    \item $K_e$ is not 0.0265 V·s/rad (datasheet error or units confusion)
    \item There is a gearbox we did not account for
    \item The voltage measurement is incorrect
\end{itemize}

\textbf{Lesson:} This example illustrates the importance of validating datasheet values. The discrepancy between predicted and measured steady-state velocity indicates a modeling error or incorrect parameter.

For a proper identification, we would:
\begin{enumerate}
    \item Verify all electrical constants experimentally (measure $R_a$, $K_t$, $K_e$ independently)
    \item Check for hidden dynamics (gearbox ratio, compliance)
    \item Use current control instead of voltage control to bypass electrical dynamics
\end{enumerate}

\section{Example 2: Thermal System Time Constant Identification}

\subsection{System Description}

We have a heating element in a metal block. We want to identify the thermal time constant $\tau$ and the heat transfer coefficient.

\textbf{Physical model:}
\begin{equation}
mc\frac{dT}{dt} = Q_{\text{in}} - hA(T - T_{\infty})
\end{equation}

where:
\begin{itemize}
    \item $m = 0.5 \, \text{kg}$ (mass of block)
    \item $c = 900 \, \text{J/(kg}\cdot\text{K)}$ (specific heat of aluminum)
    \item $h$ = unknown heat transfer coefficient
    \item $A = 0.01 \, \text{m}^2$ (surface area)
    \item $T_{\infty} = 20°\text{C}$ (ambient temperature)
    \item $Q_{\text{in}}$ = heater power (W)
\end{itemize}

Rearranging:
\begin{equation}
\frac{dT}{dt} = -\frac{hA}{mc}(T - T_{\infty}) + \frac{1}{mc}Q_{\text{in}}
\end{equation}

Define:
\begin{align}
\tau &= \frac{mc}{hA} \quad \text{(thermal time constant)} \\
K &= \frac{1}{hA} \quad \text{(gain)}
\end{align}

The model becomes:
\begin{equation}
\frac{dT}{dt} = -\frac{1}{\tau}(T - T_{\infty}) + \frac{K}{\tau}Q_{\text{in}}
\end{equation}

For a step input $Q_{\text{in}} = Q_0$:
\begin{equation}
T(t) = T_{\infty} + KQ_0(1 - e^{-t/\tau})
\end{equation}

\subsection{Experiment}

\begin{itemize}
    \item Initial temperature: $T(0) = 20°\text{C}$
    \item Heater power: $Q_0 = 50 \, \text{W}$
    \item Sampling rate: 1 Hz (1 sample per second)
    \item Duration: 600 seconds
\end{itemize}

\subsection{Data}

Temperature measurements show:
\begin{itemize}
    \item Steady-state temperature: $T_{\infty, \text{heated}} = 80°\text{C}$
    \item Temperature rise: $\Delta T = 60°\text{C}$
    \item Time to reach 63.2\% of rise: $t_{0.632} = 120 \, \text{s}$
\end{itemize}

\subsection{Parameter Estimation}

From the graphical method:
\begin{equation}
\tau = 120 \, \text{s}
\end{equation}

The gain is:
\begin{equation}
K = \frac{\Delta T}{Q_0} = \frac{60}{50} = 1.2 \, \text{K/W}
\end{equation}

Extract the heat transfer coefficient:
\begin{equation}
hA = \frac{1}{K} = \frac{1}{1.2} = 0.833 \, \text{W/K}
\end{equation}

\begin{equation}
h = \frac{0.833}{A} = \frac{0.833}{0.01} = 83.3 \, \text{W/(m}^2\cdot\text{K)}
\end{equation}

This is a reasonable value for natural convection in air.

Verify the time constant:
\begin{equation}
\tau = \frac{mc}{hA} = \frac{0.5 \times 900}{0.833} = 540 \, \text{s}
\end{equation}

\textbf{Discrepancy:} We measured $\tau = 120 \, \text{s}$ but calculated $540 \, \text{s}$ from first principles.

\subsection{Diagnosis}

Possible causes:
\begin{enumerate}
    \item The effective thermal mass is smaller than the block mass (only part of the block heats)
    \item Heat transfer is dominated by conduction to the mounting surface, not convection to air
    \item The thermocouple is not measuring the bulk temperature
\end{enumerate}

\textbf{Solution:} Insulate the block from the mounting surface and remeasure. If $\tau$ increases significantly, conduction to the mount was dominant.

\textbf{Lesson:} Discrepancies between measured and calculated parameters reveal modeling assumptions that are violated. Use them as diagnostic tools.

\section{Example 3: Spring Stiffness and Damping Measurement}

\subsection{System Description}

A mass-spring-damper system with:
\begin{equation}
m\ddot{x} + b\dot{x} + kx = F(t)
\end{equation}

\textbf{Known:} Mass $m = 1 \, \text{kg}$

\textbf{Unknown:} Spring stiffness $k$, damping coefficient $b$

\subsection{Experiment: Free Oscillation (Decay Test)}

Displace the mass by $x_0 = 0.1 \, \text{m}$ and release from rest. Measure the position $x(t)$ as it oscillates and decays.

The solution for an underdamped system is:
\begin{equation}
x(t) = Ae^{-\zeta\omega_n t}\cos(\omega_d t + \phi)
\end{equation}

where:
\begin{align}
\omega_n &= \sqrt{k/m} \quad \text{(natural frequency)} \\
\zeta &= \frac{b}{2\sqrt{mk}} \quad \text{(damping ratio)} \\
\omega_d &= \omega_n\sqrt{1 - \zeta^2} \quad \text{(damped frequency)}
\end{align}

\subsection{Data Analysis}

From the measured oscillations:
\begin{itemize}
    \item Period of oscillation: $T = 0.5 \, \text{s}$
    \item First peak: $x_1 = 0.1 \, \text{m}$
    \item Second peak: $x_2 = 0.07 \, \text{m}$ (one period later)
\end{itemize}

The damped frequency is:
\begin{equation}
\omega_d = \frac{2\pi}{T} = \frac{2\pi}{0.5} = 12.57 \, \text{rad/s}
\end{equation}

The logarithmic decrement is:
\begin{equation}
\delta = \ln\left(\frac{x_1}{x_2}\right) = \ln\left(\frac{0.1}{0.07}\right) = 0.357
\end{equation}

The damping ratio is:
\begin{equation}
\zeta = \frac{\delta}{\sqrt{4\pi^2 + \delta^2}} = \frac{0.357}{\sqrt{4\pi^2 + 0.357^2}} = 0.0567
\end{equation}

The natural frequency is:
\begin{equation}
\omega_n = \frac{\omega_d}{\sqrt{1 - \zeta^2}} = \frac{12.57}{\sqrt{1 - 0.0567^2}} = 12.59 \, \text{rad/s}
\end{equation}

Extract stiffness:
\begin{equation}
k = m\omega_n^2 = 1 \times 12.59^2 = 158.5 \, \text{N/m}
\end{equation}

Extract damping:
\begin{equation}
b = 2\zeta\sqrt{mk} = 2 \times 0.0567 \times \sqrt{1 \times 158.5} = 1.43 \, \text{N}\cdot\text{s/m}
\end{equation}

\subsection{Validation}

We apply a sinusoidal force at the resonant frequency $\omega_n$ and measure the amplitude ratio. At resonance:
\begin{equation}
\left|\frac{X}{F/k}\right| = \frac{1}{2\zeta} = \frac{1}{2 \times 0.0567} = 8.82
\end{equation}

\textbf{Experimental result:} Amplitude ratio = 9.1

This is within 3\% of the prediction, validating our identification.

\section{Example 4: Fluid System Resistance Estimation}

\subsection{System Description}

A water tank with outlet valve. The dynamics are:
\begin{equation}
A\frac{dh}{dt} = Q_{\text{in}} - Q_{\text{out}}
\end{equation}

where $h$ is the water level, $A$ is the tank cross-sectional area, $Q_{\text{in}}$ is the inlet flow rate, and $Q_{\text{out}}$ is the outlet flow rate.

For laminar flow through a valve, $Q_{\text{out}} = h/R$ where $R$ is the hydraulic resistance.

\begin{equation}
A\frac{dh}{dt} = Q_{\text{in}} - \frac{h}{R}
\end{equation}

\textbf{Known:} $A = 0.1 \, \text{m}^2$

\textbf{Unknown:} $R$ (hydraulic resistance)

\subsection{Experiment}

Fill the tank to $h_0 = 1 \, \text{m}$, turn off the inlet ($Q_{\text{in}} = 0$), and measure the water level as it drains.

The solution is:
\begin{equation}
h(t) = h_0 e^{-t/(AR)}
\end{equation}

Define the time constant $\tau = AR$.

\subsection{Data}

\begin{itemize}
    \item Initial level: $h_0 = 1 \, \text{m}$
    \item Level after 100 s: $h(100) = 0.6 \, \text{m}$
\end{itemize}

From:
\begin{equation}
h(t) = h_0 e^{-t/\tau}
\end{equation}

We have:
\begin{equation}
0.6 = 1 \times e^{-100/\tau}
\end{equation}

Take logarithm:
\begin{equation}
\ln(0.6) = -\frac{100}{\tau}
\end{equation}

\begin{equation}
\tau = -\frac{100}{\ln(0.6)} = -\frac{100}{-0.511} = 195.7 \, \text{s}
\end{equation}

Extract resistance:
\begin{equation}
R = \frac{\tau}{A} = \frac{195.7}{0.1} = 1957 \, \text{s/m}^2
\end{equation}

\subsection{Validation}

We verify by measuring the steady-state flow rate for a constant inlet $Q_{\text{in}} = 0.001 \, \text{m}^3\text{/s}$.

At steady state, $dh/dt = 0$, so:
\begin{equation}
Q_{\text{in}} = \frac{h_{\infty}}{R} \quad \Rightarrow \quad h_{\infty} = RQ_{\text{in}} = 1957 \times 0.001 = 1.957 \, \text{m}
\end{equation}

\textbf{Measured:} $h_{\infty} = 1.9 \, \text{m}$

Excellent agreement (3\% error).

\subsection{Nonlinear Extension}

For turbulent flow, $Q_{\text{out}} = C\sqrt{h}$ (Torricelli's law). The model becomes:
\begin{equation}
A\frac{dh}{dt} = Q_{\text{in}} - C\sqrt{h}
\end{equation}

This is nonlinear in $h$. To identify $C$, we linearize around a nominal operating point $h_0$:
\begin{equation}
Q_{\text{out}} \approx C\sqrt{h_0} + \frac{C}{2\sqrt{h_0}}(h - h_0)
\end{equation}

Or we fit the nonlinear model directly using nonlinear least-squares.

\section{Example 5: Mass Estimation from Force-Acceleration Data}

\subsection{System Description}

Newton's second law:
\begin{equation}
F = ma
\end{equation}

We measure applied force $F(t)$ and resulting acceleration $a(t)$ to estimate mass $m$.

\subsection{Experiment}

Apply a known force using a calibrated load cell, and measure acceleration with an accelerometer.

\textbf{Data:} 100 measurements of $(F[k], a[k])$ pairs.

\subsection{Least-Squares Estimation}

The model is:
\begin{equation}
F[k] = m \cdot a[k] + e[k]
\end{equation}

This is linear in $m$ (one parameter). The regressor is:
\begin{equation}
\mathbf{\Phi} = \begin{bmatrix} a[1] \\ a[2] \\ \vdots \\ a[100] \end{bmatrix}, \quad
\mathbf{y} = \begin{bmatrix} F[1] \\ F[2] \\ \vdots \\ F[100] \end{bmatrix}
\end{equation}

The least-squares solution is:
\begin{equation}
\hat{m} = \frac{\sum_{k=1}^{100} F[k]a[k]}{\sum_{k=1}^{100} a[k]^2}
\end{equation}

\subsection{Numerical Example}

Given:
\begin{align}
\sum F[k]a[k] &= 250 \, \text{N}\cdot\text{m/s}^2 \\
\sum a[k]^2 &= 125 \, (\text{m/s}^2)^2
\end{align}

\begin{equation}
\hat{m} = \frac{250}{125} = 2.0 \, \text{kg}
\end{equation}

\subsection{Uncertainty}

The variance of the estimate is:
\begin{equation}
\text{Var}(\hat{m}) = \frac{\sigma^2}{\sum a[k]^2}
\end{equation}

If the force measurement noise has standard deviation $\sigma = 0.5 \, \text{N}$:
\begin{equation}
\text{Var}(\hat{m}) = \frac{0.5^2}{125} = 0.002 \, \text{kg}^2
\end{equation}

\begin{equation}
\text{SE}(\hat{m}) = \sqrt{0.002} = 0.045 \, \text{kg}
\end{equation}

The 95\% confidence interval is:
\begin{equation}
\hat{m} \pm 2 \times 0.045 = 2.0 \pm 0.09 \, \text{kg}
\end{equation}

This is a 4.5\% uncertainty, which is excellent.

\section{Example 6: Coupled System Parameter Identification}

\subsection{System Description}

Two masses connected by a spring:
\begin{align}
m_1 \ddot{x}_1 &= -k(x_1 - x_2) + F_1 \\
m_2 \ddot{x}_2 &= k(x_1 - x_2) + F_2
\end{align}

\textbf{Known:} $m_1 = 1 \, \text{kg}$, $m_2 = 0.5 \, \text{kg}$

\textbf{Unknown:} Spring stiffness $k$

\subsection{Experiment}

Apply a force $F_1$ to mass 1, keep $F_2 = 0$, and measure both positions $x_1(t)$ and $x_2(t)$.

\subsection{State-Space Form}

Define state vector:
\begin{equation}
\mathbf{x} = \begin{bmatrix} x_1 \\ x_2 \\ \dot{x}_1 \\ \dot{x}_2 \end{bmatrix}
\end{equation}

The state equations are:
\begin{equation}
\dot{\mathbf{x}} = \begin{bmatrix}
0 & 0 & 1 & 0 \\
0 & 0 & 0 & 1 \\
-k/m_1 & k/m_1 & 0 & 0 \\
k/m_2 & -k/m_2 & 0 & 0
\end{bmatrix} \mathbf{x} + \begin{bmatrix} 0 \\ 0 \\ 1/m_1 \\ 0 \end{bmatrix} F_1
\end{equation}

\subsection{Identification Approach}

Discretize and write:
\begin{equation}
\mathbf{x}[k+1] = \mathbf{A}_d \mathbf{x}[k] + \mathbf{B}_d F_1[k]
\end{equation}

Since $k$ appears in $\mathbf{A}_d$, this is a nonlinear identification problem.

\textbf{Simplification:} Use only the acceleration equations:
\begin{align}
\ddot{x}_1[k] &= -\frac{k}{m_1}(x_1[k] - x_2[k]) + \frac{F_1[k]}{m_1} \\
\ddot{x}_2[k] &= \frac{k}{m_2}(x_1[k] - x_2[k])
\end{align}

Approximate accelerations using finite differences:
\begin{equation}
\ddot{x}_1[k] \approx \frac{x_1[k+1] - 2x_1[k] + x_1[k-1]}{(\Delta t)^2}
\end{equation}

Substitute:
\begin{equation}
\frac{x_1[k+1] - 2x_1[k] + x_1[k-1]}{(\Delta t)^2} = -\frac{k}{m_1}(x_1[k] - x_2[k]) + \frac{F_1[k]}{m_1}
\end{equation}

Rearrange:
\begin{equation}
x_1[k+1] - 2x_1[k] + x_1[k-1] = -\frac{k(\Delta t)^2}{m_1}(x_1[k] - x_2[k]) + \frac{(\Delta t)^2}{m_1}F_1[k]
\end{equation}

This is linear in $k$. Define $\theta = k(\Delta t)^2 / m_1$:
\begin{equation}
x_1[k+1] = (2 - \theta)x_1[k] - x_1[k-1] + \theta x_2[k] + \frac{(\Delta t)^2}{m_1}F_1[k]
\end{equation}

The regressor for mass 1 is:
\begin{equation}
x_1[k+1] = \theta_1 x_1[k] + \theta_2 x_1[k-1] + \theta_3 x_2[k] + \theta_4 F_1[k]
\end{equation}

where $\theta_1 = 2 - \theta$, $\theta_2 = -1$, $\theta_3 = \theta$, $\theta_4 = (\Delta t)^2/m_1$.

Since $\theta_2 = -1$ is known, we fit:
\begin{equation}
x_1[k+1] + x_1[k-1] = \theta_1 x_1[k] + \theta_3 x_2[k] + \theta_4 F_1[k]
\end{equation}

Least-squares gives $\hat{\theta}_1$, $\hat{\theta}_3$, $\hat{\theta}_4$.

Extract:
\begin{equation}
k = \frac{m_1 \theta_3}{(\Delta t)^2}
\end{equation}

\subsection{Validation}

Use the identified $k$ to predict $x_2(t)$ from the second equation and compare to measurements.

\section{Summary and Best Practices}

\subsection{Key Takeaways}

\begin{enumerate}
    \item \textbf{Gray-box models are preferred:} Use physical structure and estimate numerical parameters.

    \item \textbf{Experimental design is critical:} Good data beats fancy algorithms.

    \item \textbf{Least-squares is the workhorse:} Simple, robust, and well-understood.

    \item \textbf{Always validate:} Check residuals, cross-validate, verify physical plausibility.

    \item \textbf{Quantify uncertainty:} Parameter estimates without confidence intervals are incomplete.

    \item \textbf{Iterate:} Model refinement is a loop, not a one-shot process.

    \item \textbf{Trust but verify:} Datasheets are a starting point, not gospel.
\end{enumerate}

\subsection{Practical Workflow}

\begin{enumerate}
    \item Build first-principles model with nominal parameters
    \item Design experiment to excite relevant dynamics
    \item Collect high-quality data (good SNR, sufficient duration)
    \item Fit parameters using least-squares
    \item Validate on new data
    \item Check residuals for systematic errors
    \item Refine model structure if needed
    \item Quantify uncertainty
    \item Deploy and monitor
\end{enumerate}

\subsection{Common Pitfalls}

\begin{itemize}
    \item Poor signal-to-noise ratio (use larger excitation amplitudes)
    \item Insufficient excitation (use rich input signals)
    \item Validating on training data (always split train/test)
    \item Ignoring physical constraints (negative inertia is impossible)
    \item Over-parameterization (more parameters than data can support)
    \item Trusting nominal values blindly (measure what matters)
    \item Neglecting uncertainty (confidence intervals are not optional)
\end{itemize}

\subsection{When to Stop Refining}

You have a good enough model when:
\begin{itemize}
    \item Validation error is acceptable for your control objectives
    \item Residuals show no obvious patterns
    \item Parameters are physically plausible and consistent
    \item Confidence intervals are tight enough for reliable design
    \item Further refinement does not improve control performance
\end{itemize}

Engineering is about trade-offs. A simple model identified reliably often outperforms a complex model with uncertain parameters.

\section{Further Reading}

\begin{itemize}
    \item Ljung, L. \textit{System Identification: Theory for the User}. Prentice Hall, 1999. (The definitive reference)
    \item Åström, K. J., and Wittenmark, B. \textit{Adaptive Control}. Dover, 2008. (Chapter 2 on parameter estimation)
    \item Keesman, K. J. \textit{System Identification: An Introduction}. Springer, 2011. (Accessible introduction)
    \item Pintelon, R., and Schoukens, J. \textit{System Identification: A Frequency Domain Approach}. IEEE Press, 2012. (Frequency-domain methods)
\end{itemize}

\section{Problems}

\subsection{Problem 1: RC Circuit Time Constant}

An RC circuit has unknown resistance $R$ and capacitance $C$. You apply a 5 V step input and measure the voltage across the capacitor. At $t = 0.1 \, \text{s}$, the voltage is 3.16 V. Estimate the time constant $\tau = RC$.

\subsection{Problem 2: Pendulum Damping}

A pendulum has length $L = 1 \, \text{m}$ and mass $m = 0.5 \, \text{kg}$. You release it from an angle of 10° and measure the amplitude of oscillation after each swing. The first swing reaches 9°, the second 8.1°. Estimate the damping coefficient $b$.

\subsection{Problem 3: Motor Parameter Sensitivity}

For the DC motor in Example 1, assume the torque constant $K_t$ has a 10\% uncertainty. Propagate this uncertainty through the calculation to find the uncertainty in the identified inertia $J$.

\subsection{Problem 4: Frequency Response Fitting}

Given magnitude measurements $|G(j\omega)|$ at frequencies $\omega = [1, 2, 5, 10, 20] \, \text{rad/s}$ with values $[0.95, 0.89, 0.71, 0.45, 0.20]$, fit a first-order model $G(s) = K/(\tau s + 1)$ using least-squares.

\subsection{Problem 5: MIMO Spring System}

Two masses $m_1 = m_2 = 1 \, \text{kg}$ are connected by a spring of unknown stiffness $k$. You measure the natural frequencies of the coupled system as $\omega_1 = 5 \, \text{rad/s}$ and $\omega_2 = 15 \, \text{rad/s}$. Estimate $k$.

\textit{Hint: The natural frequencies of a two-mass-one-spring system are $\omega = \sqrt{k/m}$ and $\omega = \sqrt{3k/m}$.}
